\documentclass{article}
\usepackage[utf8]{inputenc}
\usepackage{amsmath,amssymb}
\usepackage{booktabs}
\usepackage{graphicx}
\usepackage{hyperref}
\usepackage{geometry}
\geometry{margin=1in}

\begin{document}

\title{The Limits of Unsupervised Absorption Detection in Sparse Autoencoders: A Systematic Analysis}
\maketitle

\subsection{Abstract}

Sparse Autoencoders (SAEs) are the dominant tool for mechanistic interpretability, but \textit{feature absorption}---where parent features suppress child features---undermines their reliability. While supervised detection methods exist, the community has sought unsupervised alternatives. This paper presents a systematic empirical investigation of the leading unsupervised approach: co-occurrence clustering via the Unsupervised Absorption Detector (UAD). Our key finding is that \textbf{UAD performs no better than random for absorption detection at scale}: on 500 active features from a GPT-2 Small SAE, UAD achieves F1 = 0.007 with 0.37\% precision, indistinguishable from random selection (F1 = 0.0075). Through controlled ablations and theoretical analysis, we show that the fundamental issue is conceptual: co-occurrence clustering detects \textit{correlation patterns}, while absorption requires detecting \textit{suppression signals}---fundamentally different statistical phenomena. We formalize this distinction and argue that absorption detection is inherently a causal inference task. Our work provides the first systematic empirical analysis of unsupervised absorption detection at scale, guiding the community toward either supervised detection or preventive architectures.


\section{1. Introduction}

Sparse Autoencoders (SAEs) have become the dominant tool for extracting interpretable features from neural network activations [Bricken et al., 2023; Cunningham et al., 2023]. The core promise is that each SAE dictionary element corresponds to a semantically coherent concept---what has been termed \textit{monosemanticity}. Yet a persistent concern undermines this promise: \textit{feature absorption}, where distinct semantic concepts are mapped to shared dictionary elements, potentially degrading the reliability of interpretability analyses.

Feature absorption was first formally characterized by Chanin et al. [2024], who demonstrated that parent features (e.g., "starts with a vowel") can suppress child features (e.g., "starts with 'a'") when both co-occur, causing the SAE to fire only the more general feature. This creates systematic "holes" in feature coverage: the SAE appears to have learned a clean monosemantic representation, but silently fails to activate for subsets of its semantic domain. While Chanin et al.'s detection protocol requires known parent-child hierarchies, the community has sought unsupervised alternatives that would enable absorption detection on any concept set without labeled supervision.

\textbf{The Central Question.} Can feature absorption be detected without ground-truth hierarchy labels? Prior work has proposed co-occurrence clustering as a potential solution: features that frequently co-occur may indicate absorption [Chen et al., 2025]. This paper subjects this hypothesis to rigorous empirical testing at scale.

\textbf{Our Contribution.} We present the first systematic empirical analysis of unsupervised absorption detection at scale. Our key finding: \textbf{UAD achieves near-random performance for absorption detection (F1 = 0.007, precision = 0.37\% on 3,702 same-cluster pairs), statistically indistinguishable from random selection (F1 = 0.0075)}. Through controlled ablations and theoretical analysis, we show that the fundamental issue is conceptual, not implementational: absorption requires detecting \textit{suppression signals} (one feature preventing another from activating), whereas co-occurrence clustering detects \textit{correlation patterns} (features that frequently appear together). These are fundamentally different statistical phenomena.

Our analysis yields three contributions:
\begin{enumerate}
  \item \textbf{Empirical Characterization}: We provide the first empirical demonstration that UAD performs no better than random for absorption detection at scale on a real SAE.
  \item \textbf{Theoretical Insight}: We formalize the distinction between correlation (what clustering detects) and suppression (what absorption requires), explaining why the former cannot proxy for the latter.
  \item \textbf{Mitigation Feasibility}: We demonstrate that inference-time mitigation (DFDA) remains feasible when absorbed pairs are known, suggesting an asymmetry: mitigation is easier than detection.
\end{enumerate}

This work reframes a promising research direction: rather than treating absorption detection as an unsupervised clustering problem, we argue it is inherently a causal inference task requiring explicit modeling of suppressive relationships.


\section{2. Related Work}

\subsection{2.1 Sparse Autoencoder Architectures}

Sparse Autoencoders project neural activations into an overcomplete sparse basis via a sparsity-inducing objective. The standard formulation uses an L1 penalty on latent activations [Makhzani and Frey, 2014]. Recent architectures have introduced alternative sparsity mechanisms: JumpReLU employs a gating mechanism that improves sparsity control [Rajamanoharan et al., 2024]; TopK enforces hard sparsity via top-k selection [Gao et al., 2024]; and Matryoshka SAEs use nested dictionaries for hierarchical feature learning [Bussmann et al., 2025]. Large-scale pre-trained SAE suites such as GemmaScope [Lieberum et al., 2024] have made cross-architecture comparison feasible. Absorption has been studied primarily in standard SAE formulations; our analysis uses GemmaScope JumpReLU SAEs.

\subsection{2.2 Feature Absorption}

Chanin et al. [2024] formally defined feature absorption as the suppression of child features by parent features under co-occurrence. Their detection protocol requires known parent-child hierarchies (e.g., first-letter spelling tasks), limiting generalization. Absorption is connected to the broader phenomenon of superposition [Elhage et al., 2022], where neural networks represent more features than neurons by using overlapping directions. Superposition creates the conditions for absorption: when multiple concepts compete for representation in overlapping directions, the SAE may suppress less dominant features in favor of more general ones.

Hierarchical SAEs (HSAE) [Chen et al., 2025] propose architectural mitigation via explicit hierarchy, but require retraining. Matryoshka SAEs [Bussmann et al., 2025] reduce collision rates from 0.49 to 0.05 through nested dictionary structure---a preventive rather than detective approach. To our knowledge, no prior work has systematically tested whether unsupervised detection of absorption is possible at scale.

\subsection{2.3 Unsupervised Detection Attempts}

Chen et al. [2025] proposed using feature co-occurrence patterns to identify related features without supervision. Co-occurrence clustering groups features that activate on similar inputs, with the intuition that absorbed features should cluster together. However, this conflates two distinct phenomena: (1) features that co-occur because they are semantically related (e.g., "cat" and "dog"), and (2) features that suppress each other (absorption). Our work is the first to empirically demonstrate that this conflation leads to near-random detection performance at scale (F1 = 0.007, indistinguishable from random F1 = 0.0075).

\subsection{2.4 Positioning}

Our work fills the gap between supervised absorption detection (Chanin et al.) and preventive architecture design (Matryoshka SAEs). We ask: if prevention is not an option, can we at least detect absorption without labels? Our negative result provides a theoretical foundation for why the field should focus on either (1) supervised detection with validated ground truth, or (2) preventive architectures that eliminate absorption at training time. The three approaches form a spectrum: supervised detection (high precision, requires labels), preventive architectures (no detection needed, requires retraining), and unsupervised detection (no labels, fails at scale). Our work rules out the third option as a viable path.


\section{3. Methodology}

\subsection{3.1 Overview}

Our experimental pipeline consists of three components: (1) the Unsupervised Absorption Detector (UAD), which uses co-occurrence clustering to flag potential absorption pairs; (2) Dynamic Feature De-Absorption (DFDA), an inference-time compensation module tested on known absorbed pairs; and (3) a random baseline to anchor F1 claims. All components are evaluated against ground-truth absorption labels obtained from Chanin et al.'s parent-child hierarchy protocol.

\subsection{3.2 Definitions and Metrics}

\textbf{Absorption.} Following Chanin et al. [2024], absorption measures parent-feature suppression of child features under co-occurrence. This requires hierarchy labels and is our gold-standard metric. Our ground-truth labels identify first-letter spelling task features that participate in known parent-child hierarchies (e.g., letters c, i, o, p, u sharing feature 18486).

\textbf{Suppression Signal.} We introduce this concept as the defining characteristic of absorption: when parent feature $p$ and child feature $c$ co-occur, the activation of $c$ is suppressed relative to its expected activation without $p$. Formally:
$$\Delta_{\text{supp}}(c, p) = \mathbb{E}[z_c \mid z_p = 0] - \mathbb{E}[z_c \mid z_p > 0]$$
where $z_i$ is the scalar activation of feature $i$ on a token. Positive $\Delta_{\text{supp}}$ indicates absorption: the child's activation is lower when the parent is present.

\subsection{3.3 The UAD Method}

We describe the Unsupervised Absorption Detector (UAD) as proposed in prior work, which we subsequently test at scale.

UAD operates in three steps:
\begin{enumerate}
  \item \textbf{Feature Filtering:} Features are filtered to the top 500 by mean activation frequency across a corpus sample of 10,000 tokens, discarding dead features (zero activation on all tokens).
  \item \textbf{Co-occurrence Matrix Construction:} For each feature $i$, compute $C_{ij} = \sum_n \mathbb{1}[z_{ni} > 0] \cdot \mathbb{1}[z_{nj} > 0]$ across the corpus sample, counting how often features $i$ and $j$ co-activate.
  \item \textbf{Hierarchical Clustering:} Cluster features by co-occurrence similarity using Ward linkage into $n_c = 50$ clusters.
  \item \textbf{Collision Detection:} Features in the same cluster with high mutual activation are flagged as potential collisions. Same-cluster pairs are all unordered pairs $(i, j)$ where features $i$ and $j$ are assigned to the same cluster.
\end{enumerate}

The intuition is that absorbed features should co-occur frequently and thus cluster together. Our experiments test whether this intuition holds at scale.

\subsection{3.4 Dynamic Feature De-Absorption (DFDA)}

DFDA is an inference-time compensation module. For each absorbed feature $i$, a small MLP (2-layer, 16 hidden units, ReLU activation, 97 parameters per pair) learns to predict the residual activation that would have been present without absorption:

$$\hat{r}_i = \text{MLP}_{\text{comp}}(z_{\text{neighbor}})$$

where $z_{\text{neighbor}}$ is a small neighborhood of co-occurring features (not the full 24,575-dimensional activation vector). The compensated reconstruction is $\hat{x}_{\text{comp}} = W_{\text{dec}}(z + \hat{r})$. Training uses MSE loss on 10,000 tokens with AdamW (learning rate $10^{-3}$) for 100 epochs.

\subsection{3.5 Random Baseline}

To anchor all F1 claims, we compute a random baseline: randomly select the same number of feature pairs as UAD detects (541 pairs), and compute F1 against known absorption labels. This baseline is essential because a non-trivial F1 can arise by chance when the detection set is large.

nd{document}
